\documentclass[10pt,a4paper]{amsart}
\usepackage[margin=19mm]{geometry}
\usepackage{amsmath,amssymb,mathtools}
\usepackage[scr=rsfs,cal=euler]{mathalfa}
\usepackage{xfrac}
\usepackage[unicode,hidelinks,bookmarks=false]{hyperref}
\newcommand{\ie}{i.e.\ }
\newcommand{\cf}{cf.\ }
\newcommand{\resp}{respectively\ }
\newcommand{\Subsection}{\subsection}
\providecommand{\nobreakdash}{\nobreak\hbox{-}}
\setlength{\emergencystretch}{2em}
\multlinegap0pt
\usepackage{bm}
\usepackage{bbm}
\usepackage{dsfont}
\usepackage{xspace}
\usepackage{cancel}
\usepackage{mathtools}
\usepackage{amsthm}
\makeatletter
\@ifundefined{theorem}{\newtheorem{theorem}{Theorem}}{}
\@ifundefined{definition}{\newtheorem{definition}[theorem]{Definition}}{}
\makeatother
\newcommand{\cdummy}{\cdot}
\newcommand{\mathd}{\mathrm{d}}
\newcommand{\tmop}[1]{\ensuremath{\operatorname{#1}}}
\title[Surface Dean--Kawasaki equations]{Surface Dean--Kawasaki equations}
\author{John Bell, Ana Djurdjevac, Nicolas Perkowski}
\date{Source: arXiv:2601.06863v2 (CC BY 4.0)}
\begin{document}
\maketitle

\noindent\emph{Source and licence.} Surface Dean--Kawasaki equations. Version arXiv:2601.06863v2 (CC BY 4.0), distributed under the Creative Commons Attribution 4.0 International licence, as recorded in the arXiv licence metadata for this exact version and shown on the abstract page https://arxiv.org/abs/2601.06863v2. Full rights notice: licenses/arxiv-2601.06863-CC-BY-4.0.txt.\par\smallskip

\noindent\textit{Editorial scope.} Counted unit: the Theorem whose source label is \texttt{None} in the article (source0.tex lines 794--805, source label \texttt{None}), delivered together with the article's own complete original proof of it (source0.tex lines 807--866, one \texttt{proof} environment of 60 lines). No mathematics is written, repaired or re-derived: statement and proof are the article's own text line for line. Required context retained uncounted: eq:eta (lines 657--660); eq:sDK-full (lines 758--769); def:sDK-martingale-problem (lines 774--786). Every definition, display and label of those lines is retained unchanged. Omitted from this package and not redistributed: 1--656, 661--757, 770--773, 787--793, 806--806, 867--1772. Nothing inside a retained excerpt is omitted or altered: every retained body is delivered exactly as the article's lines are written, so no omission receipt of any kind accompanies this package. Citations: the retained excerpts carry no citation command at all, so no bibliography entry of the article is transcribed and no reference is redistributed. Reference adaptation: none was needed and none was made; every reference inside the delivered unit resolves inside it, so no unresolved reference or citation is produced. Printed slips kept literal: no printed word, symbol, hypothesis or number of the counted statement or of its proof was altered. Layout and class: the article's own class is \texttt{article} with packages including babel, amsmath, amssymb, xcolor, amsthm, a4wide, authblk, comment, graphicx, appendix, hyperref, mathtools; the wrapper is the lane's pinned \texttt{amsart} editorial wrapper at 10pt on A4, which restates the theorem-like environments the retained text needs and adds the article's own macro definitions that the retained text needs (\texttt{\textbackslash cdummy}, \texttt{\textbackslash mathd}, \texttt{\textbackslash tmop}), with extra wrapper packages (bm, bbm, dsfont, xspace, cancel, mathtools). The delivered numbering is the unit's own. Assets: no retained span contains a figure environment, a graphics-inclusion command, a PDF-page inclusion, a table or a TikZ picture; no image file is copied into this package, the package declares no asset dependency and no image file is redistributed with it. Licence: version arXiv:2601.06863v2 of this article is distributed under the Creative Commons Attribution 4.0 International licence, as recorded in the arXiv licence metadata for this exact version and shown on the abstract page https://arxiv.org/abs/2601.06863v2. A full rights notice is delivered at licenses/arxiv-2601.06863-CC-BY-4.0.txt. Characters: every retained excerpt is pure ASCII, so the delivered engine, XeLaTeX with its default Latin Modern text font, reports no missing character.
\begin{equation}\label{eq:eta}
    \mathd \eta_t = a (\eta_t, \mu_t) \mathd t + \sigma (\eta_t, \mu_t)
   \mathd B_t,
\end{equation}

\begin{align} \nonumber
  \mathd \rho^\nu_t & = \mathd (\rho^{\nu_t}_t g_t) g_t^{- 1}\\ \nonumber
  & = g_t^{- 1} \mathd (\rho^{\nu_t}_t
  g_t) + (\rho^{\nu_t}_t g_t) \mathd g_t^{- 1} + \mathd \langle g^{- 1},
  \rho^{\nu} g \rangle_t\\ \nonumber
  & = \left( \Delta_{G_t} \rho^{\nu_t}_t + \operatorname{div}_{G_t} \left(
  \left( \nabla_{G_t} V + \int \nabla_{G_t} U (\cdummy, x) \rho^{\nu_t}_t (x)
  \nu_t (\mathd x) \right) \rho^{\nu_t}_t \right) \right) \mathd t \\
  & \quad + \sqrt{\frac2N}\operatorname{div}_{G_t} \left( \sqrt{\rho^{\nu_t}} \mathd W^G_t \right)
   - \rho^{\nu_t}_t (g_t^{- 1} \mathd g_t - g_t^{- 2} \mathd \langle g \label{eq:sDK-full}
  \rangle_t) .
\end{align}

\begin{definition}[Martingale problem]\label{def:sDK-martingale-problem}
  A stochastic process $(\mu, \eta)$ in $C (\mathbb{R}_+, \mathcal{P} (\mathcal{D})
  \times \mathbb{R}^e)$ is called a martingale solution to the
  sDK equation \eqref{eq:sDK-full}, if $\eta$ is a (probabilistically) weak  solution to \eqref{eq:eta}
  and for all $f \in C^2 (\mathcal{D})$ the process
  \[ M^f_t = \mu_t (f) - \mu_0 (f) - \int_0^t \mu_s \left( \Delta_{G_s} f -
     \nabla_{G_s} V \cdummy \nabla f - \int \nabla_{G_s} U (\cdummy, x) \mu_s
     (\mathd x) \cdummy \nabla f \right) \mathd s, \]
  $t\ge 0$, is a continuous martingale with
  \[ \langle M^f \rangle_t = \frac{2}{N} \int_0^t \mu_s (| \nabla_{G_s} f
     |_{G_s}^2) \mathd s, \qquad \langle M^f, B \rangle_t = 0, \qquad t
     \geq 0. \]
\end{definition}

\begin{theorem}[Weak uniqueness]
Let $U = 0$ and let $a$ and $\sigma$ be independent of $\mu$. Let
$\mathcal{G}_t =\mathcal{F}_t \vee \mathcal{F}^B_{\infty}$, $t\ge 0$, where
$\mathcal{F}^B$ is the canonical filtration generated by $B$. We assume that
for all $\mathcal{G}_0$-measurable $f \in C^2 (\mathcal{D})$ and $t \geq 0$ there
exists a $\mathcal{G}_0$-measurable solution $P_{\cdummy, t} f \in C^{1, 2}
(\mathbb{R}_+ \times \mathcal{D})$ to the random Kolmogorov backward equation
\[ \partial_s P_{s, t} f = - (\Delta_{G_s} P_{s, t} f - \nabla_{G_s} V \cdummy
   \nabla P_{s, t} f), \qquad P_{t, t} f = f, \]
and that strong existence and weak uniqueness hold for~\eqref{eq:eta} (which is in particular the case if $a$ and $\sigma$ are Lipschitz continuous). Then the law of the martingale solution $(\mu, \eta)$ to
the sDK Equation~\eqref{eq:sDK-full} is uniquely determined by its initial distribution.
\end{theorem}

\begin{proof}
    We start by showing that $\mathbb{E} [F (\mu_t) \mid \mathcal{G}_0]$ is uniquely
determined by $\tmop{law} (\mu_0)$ and the path $(\eta_s)_{s \geq 0}$,
for all $t \geq 0$ and $F \in C (\mathcal{P} (D))$. By a density
argument and using the linearity of $f \mapsto \mu_t (f)$, it suffices to
consider $F (\mu) = \varphi (\mu_t (f))$ for $f \in C^2$, and it is
convenient to allow $\mathcal{G}_0$-measurable random $f \in C^2$. Since
$\mu_t$ is a probability measure for each $t \geq 0$, the random variable
$\mu_t (f)$ is bounded and thus we can restrict to $\varphi (x) = x^m$ by the
Stone-Weierstra{\ss} theorem. The lemma following below this proof shows that the martingale $M^f$ in Definition~\ref{def:sDK-martingale-problem} remains a martingale in the filtration $\mathcal G$. We combine this with the telescope sum $\mu_t (P_{t, t} f) - \mu_0 (P_{0, t} f) = \lim_n \sum_{k=0}^{n-1}(\mu_{(k+1)t/n} (P_{(k+1)t/n, t} f) - \mu_{kt/n} (P_{kt/n, t} f))$ to show that for all
$\mathcal{G}_0$-measurable random $f$
\[ \mu_t (P_{t, t} f) = \mu_0 (P_{0, t} f) + \int_0^t \mu_s \left(  \left(
   \partial_s + \Delta_{G_{\eta_s}} - \nabla_{G_{\eta_s}} V \cdummy \nabla
   \right) P_{s, t} f \right) \mathd s + M^{P f}_t = \mu_0 (P_{0, t} f) + M^{P
   f}_t, \]
where $M^{P f}$ is a continuous martingale in the filtration
$(\mathcal{G}_t)_{t \geq 0}$, with quadratic variation
\[ \langle M^{P f} \rangle_s = \frac{2}{N} \int_0^s \mu_r (| \nabla_{G_r}
   P_{r, t} f |^2_{G_r}) \mathd r. \]
In particular, we obtain from It{\^o}'s formula
\begin{align} \nonumber
  & \mathbb{E} [(\mu_t (f)^m - \mu_0 (P_{0, t} f)^m) \mid \mathcal{G}_0]\\  \nonumber
  & =\mathbb{E} \left[ \left. \int_0^t m \mu_s (P_{s, t} f)^{m - 1} \mathd
  M^{P f}_s + \frac{m (m - 1)}{N} \int_0^t \mu_s (P_{s, t} f)^{m - 2} \mu_s (|
  \nabla_{G_s} P_{s, t} f |^2_{G_s}) \mathd s \,\right|\, \mathcal{G}_0 \right]\\ \label{eq:weak-uniqueness-pr1}
  & = \frac{m (m - 1)}{N} \int_0^t \mathbb{E} [\mu_s (P_{s, t} f)^{m - 2}
  \mu_s (| \nabla_{G_s} P_{s, t} f |^2_{G_s}) \mid \mathcal{G}_0] \mathd s.
\end{align}
Observe that the right hand side is $(m - 1)$-linear in $\mu_s$, which allows
an inductive proof of uniqueness: For $m = 1$, Equation~\eqref{eq:weak-uniqueness-pr1} yields
\[ \mathbb{E} [\mu_t (f) \mid \mathcal{G}_0] =\mathbb{E} [\mu_0 (P_{0, t} f)
   \mid \mathcal{G}_0], \qquad t \geq 0. \]
Assume now by induction that for all $t \geq 0$ and all
$\mathcal{G}_0$-measurable random $f_1, \ldots, f_{m - 1} \in C^2$ the
conditional expectation $\mathbb{E} [\mu_t (f_1) \cdots \mu_t (f_{m - 1})
\mid\mathcal{G}_0]$ is determined by the initial distribution of $\mu$. Then \eqref{eq:weak-uniqueness-pr1} shows that also $\mathbb{E} [\mu_t (f)^m
\mid\mathcal{G}_0]$ is uniquely determined. By
polarization, also $\mathbb{E} [\mu_t (f_1) \cdots \mu_t (f_m)
\mid\mathcal{G}_0]$ is uniquely determined for all $t \geq 0$, $m \in
\mathbb{N}$, and $\mathcal{G}_0$-measurable random $f_1, \ldots, f_m \in C^2$, which concludes the induction step.

By the tower property together with a monotone class argument, we then get
uniqueness of $\mathbb{E} [\Phi (\mu_t, \eta_t)]$ for all bounded and
measurable $\Phi : \mathcal{P} (E) \times \mathbb{R}^e \rightarrow
\mathbb{R}$. Now we perform another induction over the set of times $t_1,
\ldots, t_n$ to show that
\[ \mathbb{E} [\Phi_1 (\mu_{t_1}, \eta_{t_1}) \cdots \Phi_n (\mu_{t_n},
   \eta_{t_n})] \]
is uniquely determined by the initial distribution. We have just shown the
case $n = 1$. For general $n$, we condition on $\mathcal{F}_{t_{n - 1}}$ and
apply the previous argument to show that $\mathbb{E} [\Phi (\mu_{t_n},
\eta_{t_n}) \mid\mathcal{F}_{t_{n - 1}}]$ is uniquely determined by $ (\mu_{t_{n
- 1}}, \eta_{t_{n - 1}})$, i.e.
\[ \mathbb{E} [\Phi_1 (\mu_{t_1}, \eta_{t_1}) \cdots \Phi_{n - 1} (\mu_{t_{n -
   1}}, \eta_{t_{n - 1}}) \mathbb{E} [\Phi_n (\mu_{t_n}, \eta_{t_n})
   \mid\mathcal{F}_{t_{n - 1}}]] =\mathbb{E} [\Phi_1 (\mu_{t_1}, \eta_{t_1})
   \cdots \tilde{\Phi}_{n - 1} (\mu_{t_{n - 1}}, \eta_{t_{n - 1}})], \]
so the claim for $n$ follows from that for $n-1$. We conclude by another monotone class argument
to get the uniqueness of finite-dimensional distributions.
\end{proof}

\end{document}
